\section{Ontology Design}
\label{sec:ontology}

\subsection{Design principles}

The ontology, \emph{DrinkOnto}, was designed with three principles in mind.

\textbf{Reuse before inventing.} Standard vocabularies already cover most of
what a drinks domain needs (beverages, organisations, places, geographic
coordinates, provenance), so only genuinely domain-specific terms receive the
project namespace (\texttt{drink:}). This keeps the graph interoperable: a
query written against schema.org or Dublin Core terms matches our data directly.

\textbf{Separate types from instances.} \texttt{drink:Gin} is a \emph{class};
the individual \texttt{.../id/spirit/gin} is an instance of that class. Using
one URI as both a class and an individual (OWL punning) is avoided because
cocktails point at the individual, not at the class \cite{owl2}.

\textbf{Model the application's questions.} The schema exists to answer
questions such as ``which IBA cocktails use gin and lemon?'' and ``how much
lemon is in a Whiskey Sour?''; the second question is what motivated the
measure model described below.

\subsection{Reused vocabularies}

Table~\ref{tab:reuse} lists the external vocabularies imported by the ontology
and the terms actually used.

\begin{table}[h]
  \centering
  \small
  \begin{tabular}{@{}llp{7.2cm}@{}}
    \toprule
    Prefix & Namespace & Terms used in this project \\
    \midrule
    \texttt{schema} & schema.org & \texttt{Beverage}, \texttt{Brand},
      \texttt{Organization}, \texttt{Product}, \texttt{Place},
      \texttt{name}, \texttt{image}, \texttt{description} \\
    \texttt{dcterms} & Dublin Core Terms & \texttt{title},
      \texttt{description}, \texttt{identifier}, \texttt{source},
      \texttt{license}, \texttt{spatial} \\
    \texttt{foaf} & FOAF & \texttt{homepage}, \texttt{depiction} \\
    \texttt{geo} & WGS84 geo & \texttt{lat}, \texttt{long} \\
    \texttt{skos} & SKOS & \texttt{Concept}, \texttt{altLabel} \\
    \texttt{owl}, \texttt{rdf}, \texttt{rdfs} & W3C core & typing, labels,
      \texttt{subClassOf}, \texttt{inverseOf}, \texttt{propertyChainAxiom} \\
    \bottomrule
  \end{tabular}
  \caption{External vocabularies reused by \emph{DrinkOnto}.}
  \label{tab:reuse}
\end{table}

\subsection{Classes and properties}

Table~\ref{tab:classes} lists the classes. The umbrella class
\texttt{drink:Beverage} is a subclass of \texttt{schema:Beverage} so that
cocktails and spirits share a superclass without asserting anything about
non-alcoholic drinks.

\begin{table}[h]
  \centering
  \small
  \begin{tabular}{@{}p{4.0cm}p{3.9cm}p{6.6cm}@{}}
    \toprule
    Class & Superclass & Role \\
    \midrule
    \texttt{Cocktail} & \texttt{Beverage} & a mixed drink defined by a recipe \\
    \texttt{Spirit} & \texttt{Beverage} & a spirit type; instances are categories, not bottles \\
    \texttt{Gin}, \texttt{Whisky}, \texttt{Rum}, \texttt{Vodka}, \texttt{Tequila}, \texttt{Brandy}, \texttt{Liqueur}, \texttt{Absinthe} & \texttt{Spirit} & the eight spirit types recognised by the base-spirit classifier (Section~\ref{sec:transformation}) \\
    \texttt{Ingredient} & \texttt{schema:Product} & a component used in a recipe (juice, syrup, bitters, \dots) \\
    \texttt{IngredientAmount} & n/a & one ingredient's use in one cocktail, carrying its measure \\
    \texttt{Brand} & \texttt{schema:Brand} & a commercial product line, e.g.\ Bombay Sapphire \\
    \texttt{Distillery} & \texttt{schema:Organization} & a physical production site \\
    \texttt{GlassType} & \texttt{schema:Product} & a vessel a cocktail is served in \\
    \texttt{FlavorProfile} & \texttt{skos:Concept} & a taste or style descriptor \\
    \texttt{Beverage} & \texttt{schema:Beverage} & local umbrella class for cocktails and spirits \\
    \bottomrule
  \end{tabular}
  \caption{Classes defined in \emph{DrinkOnto}.}
  \label{tab:classes}
\end{table}

Table~\ref{tab:objectprops} lists the object properties. Datatype properties
attach literal values: \texttt{isIBA} (boolean), \texttt{preparation},
\texttt{garnish} and \texttt{hasTag} (strings) on cocktails;
\texttt{amountText} and \texttt{amountUnit} (strings) and \texttt{amountValue}
(decimal) on ingredient amounts; \texttt{foundedYear} (gYear) on distilleries;
and \texttt{aboutABV} (decimal) on spirits.

\begin{table}[h]
  \centering
  \small
  \begin{tabular}{@{}llp{5.6cm}@{}}
    \toprule
    Property & Domain $\rightarrow$ Range & Notes \\
    \midrule
    \texttt{baseSpirit} & Cocktail $\rightarrow$ Spirit & the dominant spirit \\
    \texttt{usesIngredient} & Cocktail $\rightarrow$ IngredientAmount & recipe edge carrying a measure \\
    \texttt{ofIngredient} & IngredientAmount $\rightarrow$ Ingredient & the ingredient a usage refers to \\
    \texttt{hasIngredient} & Cocktail $\rightarrow$ Ingredient & not asserted: derived by the property chain \\
    \texttt{servedIn} & Cocktail $\rightarrow$ GlassType & \\
    \texttt{producedBy} / \texttt{produces} & Brand $\leftrightarrow$ Distillery & inverse pair \\
    \texttt{brandOf} / \texttt{hasBrand} & Brand $\leftrightarrow$ Spirit & inverse pair \\
    \texttt{locatedIn} & Distillery $\rightarrow$ \texttt{schema:Place} & sub-property of \texttt{dcterms:spatial} and \texttt{schema:location} \\
    \texttt{countryOfOrigin} & Distillery $\rightarrow$ \texttt{schema:Country} & \\
    \texttt{hasFlavorProfile} & Beverage $\rightarrow$ FlavorProfile & \\
    \bottomrule
  \end{tabular}
  \caption{Object properties: domain, range and design notes.}
  \label{tab:objectprops}
\end{table}

\subsection{The measure model}

The hardest modelling question was ``how much lemon is in a Whiskey Sour?''. A
plain \texttt{Cocktail} $\rightarrow$ \texttt{Ingredient} edge cannot carry a
measure, because an RDF predicate has no fields of its own. The solution is to
make each ingredient \emph{usage} a first-class node of type
\texttt{drink:IngredientAmount}. Listing~\ref{lst:amount} shows the pattern for
the Negroni.

\begin{lstlisting}[frame=single, basicstyle=\ttfamily\small,
    caption={Ingredient-amount pattern for the Negroni.},
    label={lst:amount}]
id:ingredient-amount/negroni/1  a drink:IngredientAmount ;
        drink:ofIngredient  id:ingredient/gin ;
        drink:amountText    "1 oz" ;
        drink:amountValue   1.0 ;
        drink:amountUnit    "oz" .

id:cocktail/negroni  drink:usesIngredient  id:ingredient-amount/negroni/1 .
\end{lstlisting}

The measure is now addressable: every quantity can be queried or compared, and
the prose form of the source is preserved next to the parsed numeric value.
Section~\ref{sec:transformation} describes how the numeric value is extracted.

Finally, the ontology declares an OWL~2 property chain so that
\texttt{hasIngredient} does not have to be stored at all:

\begin{lstlisting}[frame=single, basicstyle=\ttfamily\small, breaklines=true]
drink:hasIngredient  owl:propertyChainAxiom  ( drink:usesIngredient drink:ofIngredient )
\end{lstlisting}

A reasoner could use this axiom to derive \texttt{negroni hasIngredient gin}.
The pipeline does not run a reasoner; instead, the demonstration queries
traverse the same chain with a SPARQL property path
(\texttt{usesIngredient/ofIngredient}), which yields identical
answers without materialising inferred triples.

\subsection{The ontology at a glance}

Figure~\ref{fig:ontology} is the VOWL rendering of DrinkOnto exported from
WebVOWL. Classes are circles. External classes, including those reused from
schema.org and SKOS and the upper RDFS and OWL vocabulary, are the darker
circles marked ``(external)''. Object properties are blue rectangles and
datatype properties are green; yellow rectangles are datatypes. Dotted arrows
are \texttt{rdfs:subClassOf}. An inverse pair shares one edge and carries both
labels.

\begin{figure}[p]
  \centering
   %%%%%%%%%%%%%%%%%%%%%%%%%%%%%%%%%%%%%%%%%%%%%%%%%%%%%%%%%%%%%%%%%%%%
 %        Generated with the experimental alpha version of the TeX exporter of WebVOWL (version 1.1.3) %%% 
 %%%%%%%%%%%%%%%%%%%%%%%%%%%%%%%%%%%%%%%%%%%%%%%%%%%%%%%%%%%%%%%%%%%%

 %   The content can be used as import in other TeX documents. 
 %   Parent document has to use the following packages   
 %   \usepackage{tikz}  
 %   \usepackage{helvet}  
 %   \usetikzlibrary{decorations.markings,decorations.shapes,decorations,arrows,automata,backgrounds,petri,shapes.geometric}  
 %   \usepackage{xcolor}  

 %%%%%%%%%%%%%%% Example Parent Document %%%%%%%%%%%%%%%%%%%%%%%
 %\documentclass{article} 
 %\usepackage{tikz} 
 %\usepackage{helvet} 
 %\usetikzlibrary{decorations.markings,decorations.shapes,decorations,arrows,automata,backgrounds,petri,shapes.geometric} 
 %\usepackage{xcolor} 

 %\begin{document} 
 %\section{Example} 
 %  This is an example. 
 %  \begin{figure} 
 %    \input{<THIS_FILE_NAME>} % << tex file name for the graph 
 %    \caption{A generated graph with TKIZ using alpha version of the TeX exporter of WebVOWL (version 1.1.3) } 
 %  \end{figure} 
 %\end{document} 
 %%%%%%%%%%%%%%%%%%%%%%%%%%%%%%%%%%%%%%%%%%%%%%%%%%%%%%%%%%%%%%%%%%%%

\definecolor{imageBGCOLOR}{HTML}{FFFFFF} 
\definecolor{owlClassColor}{HTML}{AACCFF}
\definecolor{owlObjectPropertyColor}{HTML}{AACCFF}
\definecolor{owlExternalClassColor}{HTML}{AACCFF}
\definecolor{owlDatatypePropertyColor}{HTML}{99CC66}
\definecolor{owlDatatypeColor}{HTML}{FFCC33}
\definecolor{owlThingColor}{HTML}{FFFFFF}
\definecolor{valuesFrom}{HTML}{6699CC}
\definecolor{rdfPropertyColor}{HTML}{CC99CC}
\definecolor{unionColor}{HTML}{6699cc}
\begin{center} 
\resizebox{\linewidth}{!}{
\begin{tikzpicture}[framed]
\clip (-420pt , 680pt ) rectangle (1620pt , -1720pt);
\tikzstyle{dashed}=[dash pattern=on 4pt off 4pt] 
\tikzstyle{dotted}=[dash pattern=on 2pt off 2pt] 
\sffamily\fontsize{12}{12}\selectfont
 
\tikzset{triangleBlack/.style = {fill=black, draw=black, line width=1pt,scale=0.7,regular polygon, regular polygon sides=3} }
\tikzset{triangleWhite/.style = {fill=white, draw=black, line width=1pt,scale=0.7,regular polygon, regular polygon sides=3} }
\tikzset{triangleBlue/.style  = {fill=valuesFrom, draw=valuesFrom, line width=1pt,scale=0.7,regular polygon, regular polygon sides=3} }
\tikzset{Diamond/.style = {fill=white, draw=black, line width=2pt,scale=1.2,regular polygon, regular polygon sides=4} }
\tikzset{Literal/.style={rectangle,align=center,
font={\fontsize{12pt}{12}\selectfont \sffamily },
black, draw=black, dashed, line width=1pt, fill=owlDatatypeColor, minimum width=80pt,
minimum height = 20pt}}

\tikzset{Datatype/.style={rectangle,align=center,
font={\fontsize{12pt}{12}\selectfont \sffamily },
black, draw=black, line width=1pt, fill=owlDatatypeColor, minimum width=80pt,
minimum height = 20pt}}

\tikzset{owlClass/.style={circle, inner sep=0mm,align=center, 
font={\fontsize{12pt}{12}\selectfont \sffamily },
black, draw=black, line width=1pt, fill=owlClassColor, minimum size=101pt}}

\tikzset{anonymousClass/.style={circle, inner sep=0mm,align=center, 
font={\fontsize{12pt}{12}\selectfont \sffamily },
black, dashed, draw=black, line width=1pt, fill=owlClassColor, minimum size=101pt}}

\tikzset{owlThing/.style={circle, inner sep=0mm,align=center,
font={\fontsize{12pt}{12}\selectfont \sffamily },
black, dashed, draw=black, line width=1pt, fill=owlThingColor, minimum size=62pt}}

\tikzset{owlObjectProperty/.style={rectangle,align=center,
inner sep=0mm,
font={\fontsize{12pt}{12}\selectfont \sffamily },
fill=owlObjectPropertyColor, minimum width=80pt,
minimum height = 25pt}}

\tikzset{rdfProperty/.style={rectangle,align=center,
inner sep=0mm,
font={\fontsize{12pt}{12}\selectfont \sffamily },
fill=rdfPropertyColor, minimum width=80pt,
minimum height = 25pt}}

\tikzset{owlDatatypeProperty/.style={rectangle,align=center,
fill=owlDatatypePropertyColor, minimum width=80pt,
inner sep=0mm,
font={\fontsize{12pt}{12}\selectfont \sffamily },
minimum height = 25pt}}

\tikzset{rdfsSubClassOf/.style={rectangle,align=center,
font={\fontsize{12pt}{12}\selectfont \sffamily },
inner sep=0mm,
fill=imageBGCOLOR, minimum width=80pt,
minimum height = 25pt}}

\tikzset{unionOf/.style={circle, inner sep=0mm,align=center,
font={\fontsize{12pt}{12}\selectfont \sffamily },
black, draw=black, line width=1pt, fill=unionColor, minimum size=25pt}}

\tikzset{disjointWith/.style={circle, inner sep=0mm,align=center,
font={\fontsize{12pt}{12}\selectfont \sffamily },
black, draw=black, line width=1pt, fill=unionColor, minimum size=20pt}}

\tikzset{owlEquivalentClass/.style={circle,align=center,
font={\fontsize{12pt}{12}\selectfont \sffamily },
inner sep=0mm,
black, solid, draw=black, line width=3pt, fill=owlExternalClassColor, minimum size=101pt,
postaction = {draw,line width=1pt, white}}}

\draw [black,line width=2pt] plot [smooth] coordinates {(-27.516435456802498pt, -856.3385554826848pt) (-62.801927703246385pt, -941.4357439879175pt)  (-98.08741994969027pt, -1026.5329324931502pt)};
\node[triangleBlack, rotate=-202.52171053237032] at (-95.69522857666016pt, -1020.7637939453125pt)   (single_marker0) {};
 \draw [black,line width=2pt] plot [smooth] coordinates {(345.66361013271984pt, -685.8642762480888pt) (200.1260855573043pt, -733.1642104350703pt)  (54.58856098188875pt, -780.464144622052pt)};
\node[triangleBlack, rotate=-251.99605015565575] at (60.35353469848633pt, -778.5905151367188pt)   (single_marker1) {};
 \draw [black,line width=2pt] plot [smooth] coordinates {(643.6809991132841pt, -241.5384144764708pt) (586.4477910465039pt, -157.45612831345875pt)  (529.2145829797236pt, -73.37384215044673pt)};
\node[triangleBlack, rotate=34.242498279129606] at (532.8301391601562pt, -78.68545532226562pt)   (single_marker2) {};
 \draw [black,line width=2pt] plot [smooth] coordinates {(439.4874925526809pt, -616.4206142361193pt) (541.3102933016846pt, -470.95031008943954pt)  (643.1330940506884pt, -325.4800059427599pt)};
\node[triangleBlack, rotate=-34.99058774498242] at (639.617431640625pt, -330.50274658203125pt)   (single_marker3) {};
 \draw [black,line width=2pt] plot [smooth] coordinates {(1140.7426510361151pt, -217.96736274866657pt) (1259.3972058151094pt, -172.88062634840176pt)  (1378.0517605941036pt, -127.79388994813694pt)};
\node[triangleBlack, rotate=-69.19387096068769] at (1371.635498046875pt, -130.23196411132812pt)   (single_marker4) {};
 \node[triangleBlack, rotate=-249.1938709606877] at (1146.3514404296875pt, -215.8361358642578pt)   (INV_single_marker4) {};
 \draw [black, dotted ,line width=2pt] plot [smooth] coordinates {(779.5483574703727pt, -1554.4866308301891pt) (883.2053582945794pt, -1489.7772654988203pt)  (986.8623591187861pt, -1425.0679001674512pt)};
\node[triangleWhite, rotate=-58.02547657468389] at (981.4386596679688pt, -1428.4537353515625pt)   (single_marker5) {};
 \draw [black, dotted ,line width=2pt] plot [smooth] coordinates {(-278.64717831611426pt, -481.81154160272973pt) (-291.33494755460697pt, -357.49574569420247pt)  (-304.0227167930997pt, -233.17994978567526pt)};
\node[triangleWhite, rotate=5.8278152709933835] at (-303.31976318359375pt, -240.0673370361328pt)   (single_marker6) {};
 \draw [black, dotted ,line width=2pt] plot [smooth] coordinates {(734.8983423166935pt, -1359.665038357244pt) (857.224418122161pt, -1375.5742224119872pt)  (979.5504939276284pt, -1391.4834064667307pt)};
\node[triangleWhite, rotate=-97.40988581868466] at (972.8939819335938pt, -1390.61767578125pt)   (single_marker7) {};
 \draw [black, dotted ,line width=2pt] plot [smooth] coordinates {(868.6273081440115pt, 546.9477881374667pt) (986.630409881819pt, 526.1067937799369pt)  (1104.6335116196265pt, 505.26579942240716pt)};
\node[triangleWhite, rotate=-100.01610818049413] at (1098.0762939453125pt, 506.4239196777344pt)   (single_marker8) {};
 \draw [black,line width=2pt] plot [smooth] coordinates {(830.3327671636933pt, -648.8638575230209pt) (973.4832324960514pt, -677.6298083669317pt)  (1116.6336978284094pt, -706.3957592108425pt)};
\node[triangleBlack, rotate=-101.3618837697186] at (1110.7275390625pt, -705.2089233398438pt)   (single_marker9) {};
 \draw [black, dotted ,line width=2pt] plot [smooth] coordinates {(1208.273750290044pt, -686.9959800618542pt) (1316.7622232849749pt, -610.2734221110867pt)  (1425.2506962799057pt, -533.5508641603192pt)};
\node[triangleWhite, rotate=-54.7326619339291] at (1419.737548828125pt, -537.44970703125pt)   (single_marker10) {};
 \draw [black, dotted ,line width=2pt] plot [smooth] coordinates {(625.2736128852645pt, 40.17535451989511pt) (645.1558698859159pt, -96.52697974045292pt)  (665.0381268865673pt, -233.22931400080094pt)};
\node[triangleWhite, rotate=-171.72434312502085] at (664.1340942382812pt, -227.01345825195312pt)   (single_marker11) {};
 \draw [black, dotted ,line width=2pt] plot [smooth] coordinates {(805.9584742213871pt, -682.9105183626029pt) (872.5057568819805pt, -797.4164753680756pt)  (939.0530395425739pt, -911.9224323735482pt)};
\node[triangleWhite, rotate=-149.83632304754758] at (935.5968017578125pt, -905.975341796875pt)   (single_marker12) {};
 \draw [black, dotted ,line width=2pt] plot [smooth] coordinates {(419.2995289639969pt, -329.39215049463064pt) (327.497626294213pt, -226.992965142877pt)  (235.69572362442906pt, -124.59377979112338pt)};
\node[triangleWhite, rotate=41.876472237863425] at (239.7345428466797pt, -129.0988311767578pt)   (single_marker13) {};
 \draw [black, dotted ,line width=2pt] plot [smooth] coordinates {(1140.8833217693837pt, -261.31570117548756pt) (1267.7888631777373pt, -308.5974100220585pt)  (1394.6944045860912pt, -355.87911886862946pt)};
\node[triangleWhite, rotate=-110.43407363166438] at (1388.27099609375pt, -353.48590087890625pt)   (single_marker14) {};
 \draw [black, dotted ,line width=2pt] plot [smooth] coordinates {(465.5701969078026pt, -662.2685169415001pt) (597.5217350357291pt, -652.4371088205667pt)  (729.4732731636556pt, -642.6057006996332pt)};
\node[triangleWhite, rotate=-85.73928960598847] at (722.8569946289062pt, -643.0986938476562pt)   (single_marker15) {};
 \draw [black, dotted ,line width=2pt] plot [smooth] coordinates {(687.2118930418696pt, -332.4934826230176pt) (726.355368355004pt, -461.2573117389028pt)  (765.4988436681384pt, -590.021140854788pt)};
\node[triangleWhite, rotate=-163.09106947205066] at (763.7059936523438pt, -584.1234741210938pt)   (single_marker16) {};
 \draw [black, dotted ,line width=2pt] plot [smooth] coordinates {(1398.3704401748594pt, -53.334970967081034pt) (1318.9305452273738pt, 45.939166406876pt)  (1239.490650279888pt, 145.21330378083303pt)};
\node[triangleWhite, rotate=38.666635129324476] at (1243.421630859375pt, 140.30087280273438pt)   (single_marker17) {};
 \draw [black,line width=2pt] plot [smooth] coordinates {(723.0901351622915pt, -278.28317408229395pt) (872.5518902451008pt, -262.3232245218422pt)  (1022.01364532791pt, -246.36327496139037pt)};
\node[triangleBlack, rotate=-83.9049866701524] at (1015.4281616210938pt, -247.06649780273438pt)   (single_marker18) {};
 \draw [black,line width=2pt] plot [smooth] coordinates {(349.1236189952506pt, -639.3653303944351pt) (272.9803116825276pt, -601.3761639962884pt)  (196.8370043698046pt, -563.3869975981416pt)};
\node[triangleBlack, rotate=63.48440812366047] at (202.3740234375pt, -566.1495361328125pt)   (single_marker19) {};
 \draw [black,line width=2pt] plot [smooth] coordinates {(-47.67133710887853pt, -756.1783417713939pt) (-142.45219271259987pt, -662.3072233662754pt)  (-237.23304831632123pt, -568.436104961157pt)};
\node[triangleBlack, rotate=45.276335115740125] at (-232.40350341796875pt, -573.2192993164062pt)   (single_marker20) {};
 \draw [black,line width=2pt] plot [smooth] coordinates {(398.3858734730116pt, -728.1118269049869pt) (388.8890249312997pt, -828.2751456420386pt)  (379.3921763895877pt, -928.4384643790904pt)};
\node[triangleBlack, rotate=-185.41629430868295] at (379.9797668457031pt, -922.2412109375pt)   (single_marker21) {};
 \draw [black,line width=2pt] plot [smooth] coordinates {(741.0916808540958pt, -606.2406468160071pt) (616.837873685871pt, -503.0911442471146pt)  (492.58406651764625pt, -399.941641678222pt)};
\node[triangleBlack, rotate=50.3018339868606] at (497.9537353515625pt, -404.3992919921875pt)   (single_marker22) {};
 \draw [black,line width=2pt] plot [smooth] coordinates {(1443.8338644258292pt, -169.90172796041566pt) (1453.6070119759615pt, -311.56310157930125pt)  (1463.380159526094pt, -453.2244751981868pt)};
\node[triangleBlack, rotate=-176.05394157669934] at (1462.8985595703125pt, -446.244873046875pt)   (single_marker23) {};
 \definecolor{Node24_COLOR}{HTML}{3366cc} 
 \node[owlClass ,minimum size=100pt , fill=Node24_COLOR  , text=white] at (201.65180435721177pt, -86.61995247641627pt)   (Node24) {Concept\\ {\small (external) }};
\node[owlClass ,minimum size=100pt  , text=black] at (780.3322991405279pt, -638.8163106848914pt)   (Node25) {Beverage};
\node[owlClass ,minimum size=121.09035488895913pt  , text=black] at (616.4155594837453pt, 101.07973660040156pt)   (Node26) {Tequila\\ {\color{gray} 1 }};
\definecolor{Node27_COLOR}{HTML}{3366cc} 
 \node[owlClass ,minimum size=100pt , fill=Node27_COLOR  , text=white] at (736.2861453499416pt, -1581.4936855094174pt)   (Node27) {DatatypeProperty\\ {\small (external) }};
\definecolor{Node28_COLOR}{HTML}{3366cc} 
 \node[owlClass ,minimum size=100pt , fill=Node28_COLOR  , text=white] at (1207.6261683785476pt, 185.0335849596386pt)   (Node28) {Organization\\ {\small (external) }};
\node[owlClass ,minimum size=100pt  , text=black] at (672.3784375694802pt, -283.6983127929141pt)   (Node29) {Spirit};
\node[owlClass ,minimum size=100pt  , text=black] at (453.3434482312142pt, -367.3659778093378pt)   (Node30) {Flavor profile};
\definecolor{Node31_COLOR}{HTML}{3366cc} 
 \node[owlClass ,minimum size=100pt , fill=Node31_COLOR  , text=white] at (964.679214623433pt, -956.0166400512596pt)   (Node31) {Beverage\\ {\small (external) }};
\node[owlClass ,minimum size=100pt  , text=black] at (-273.46897691546695pt, -532.5479804830289pt)   (Node32) {Ingredient};
\definecolor{Node33_COLOR}{HTML}{3366cc} 
 \node[owlClass ,minimum size=100pt , fill=Node33_COLOR  , text=white] at (1166.6341658515748pt, -716.443306048972pt)   (Node33) {Country\\ {\small (external) }};
\node[Datatype ,minimum width=49pt  , text=black] at (378.3496497003637pt, -939.4339996459714pt)   (Node34) {string};
\node[Datatype ,minimum width=61pt  , text=black] at (-102.6169402465575pt, -1037.4566696440647pt)   (Node35) {decimal};
\definecolor{Node36_COLOR}{HTML}{3366cc} 
 \node[owlClass ,minimum size=100pt , fill=Node36_COLOR  , text=white] at (684.3242650051047pt, -1353.0875993357517pt)   (Node36) {ObjectProperty\\ {\small (external) }};
\definecolor{Node37_COLOR}{HTML}{3366cc} 
 \node[owlClass ,minimum size=100pt , fill=Node37_COLOR  , text=white] at (1442.485201424855pt, -373.6847285890519pt)   (Node37) {Brand\\ {\small (external) }};
\node[owlClass ,minimum size=121.09035488895913pt  , text=black] at (404.1951424293711pt, -666.8414322525516pt)   (Node38) {Cocktail\\ {\color{gray} 1 }};
\node[Datatype ,minimum width=61pt  , text=black] at (521.8450788883339pt, -62.54717747757529pt)   (Node39) {decimal};
\node[Datatype ,minimum width=63pt  , text=black] at (175.8987630410628pt, -552.9405596642694pt)   (Node40) {boolean};
\node[owlClass ,minimum size=121.09035488895913pt  , text=black] at (-3.942971314762543pt, -799.4869886175892pt)   (Node41) {Ingredient amount\\ {\color{gray} 1 }};
\node[owlClass ,minimum size=121.09035488895913pt  , text=black] at (1083.2109085996817pt, -239.82845784704548pt)   (Node42) {Brand\\ {\color{gray} 1 }};
\definecolor{Node43_COLOR}{HTML}{3366cc} 
 \node[owlClass ,minimum size=100pt , fill=Node43_COLOR  , text=white] at (1154.856240383321pt, 496.39576472527216pt)   (Node43) {Class\\ {\small (external) }};
\node[owlEquivalentClass ,minimum size=129.09035488895913pt  , text=black] at (1439.3226582219138pt, -104.51197872696109pt)   (Node44) {Distillery,\\ {\small Distillery }\\ {\color{gray} 1 }};
\definecolor{Node45_COLOR}{HTML}{3366cc} 
 \node[owlClass ,minimum size=100pt , fill=Node45_COLOR  , text=white] at (1466.8902807183752pt, -504.1035381732014pt)   (Node45) {Place\\ {\small (external) }};
\definecolor{Node46_COLOR}{HTML}{3366cc} 
 \node[owlClass ,minimum size=100pt , fill=Node46_COLOR  , text=white] at (-309.200918193747pt, -182.44351090537606pt)   (Node46) {Product\\ {\small (external) }};
\definecolor{Node47_COLOR}{HTML}{3366cc} 
 \node[owlClass ,minimum size=100pt , fill=Node47_COLOR  , text=white] at (1030.1245712392172pt, -1398.0608454882229pt)   (Node47) {Property\\ {\small (external) }};
\definecolor{Node48_COLOR}{HTML}{3366cc} 
 \node[owlClass ,minimum size=100pt , fill=Node48_COLOR  , text=white] at (818.4045793803172pt, 555.8178228346017pt)   (Node48) {Class\\ {\small (external) }};
\node[owlDatatypeProperty ,minimum width=90pt  , text=black] at (-62.801927703246385pt, -941.4357439879175pt)   (property0) {amountValue};
\node[owlObjectProperty ,minimum width=99pt  , text=black] at (200.1260855573043pt, -733.1642104350703pt)   (property1) {usesIngredient};
\node[owlDatatypeProperty ,minimum width=74pt  , text=black] at (586.4477910465039pt, -157.45612831345878pt)   (property2) {aboutABV};
\node[owlObjectProperty ,minimum width=73pt  , text=black] at (541.3102933016846pt, -470.95031008943954pt)   (property3) {baseSpirit};
% Createing Inverse Property 
\node[owlObjectProperty ,minimum width=80pt  , text=black] at (1259.3972058151094pt, -186.88062634840176pt)   (property4) {produces};
% owlObjectProperty vs owlObjectProperty
% ,minimum width=80pt vs ,minimum width=84pt
%  vs 
% , text=black vs , text=black
\node[owlObjectProperty ,minimum width=84pt  , text=black] at (1259.3972058151094pt, -158.88062634840176pt)   (property4inv) {producedBy};
\node[rdfsSubClassOf ,minimum width=82pt  , text=black] at (883.2053582945794pt, -1489.7772654988203pt)   (property5) {Subclass of};
\node[rdfsSubClassOf ,minimum width=82pt  , text=black] at (-291.33494755460697pt, -357.49574569420247pt)   (property6) {Subclass of};
\node[rdfsSubClassOf ,minimum width=82pt  , text=black] at (857.224418122161pt, -1375.5742224119872pt)   (property7) {Subclass of};
\node[rdfsSubClassOf ,minimum width=82pt  , text=black] at (986.630409881819pt, 526.1067937799369pt)   (property8) {Subclass of};
\node[owlObjectProperty ,minimum width=92pt  , text=black] at (973.4832324960514pt, -677.6298083669317pt)   (property9) {origin country};
\node[rdfsSubClassOf ,minimum width=82pt  , text=black] at (1316.7622232849749pt, -610.2734221110867pt)   (property10) {Subclass of};
\node[rdfsSubClassOf ,minimum width=82pt  , text=black] at (645.155869885916pt, -96.52697974045294pt)   (property11) {Subclass of};
\node[rdfsSubClassOf ,minimum width=82pt  , text=black] at (872.5057568819805pt, -797.4164753680755pt)   (property12) {Subclass of};
\node[rdfsSubClassOf ,minimum width=82pt  , text=black] at (327.497626294213pt, -226.992965142877pt)   (property13) {Subclass of};
\node[rdfsSubClassOf ,minimum width=82pt  , text=black] at (1267.7888631777373pt, -308.5974100220585pt)   (property14) {Subclass of};
\node[rdfsSubClassOf ,minimum width=82pt  , text=black] at (597.5217350357291pt, -652.4371088205667pt)   (property15) {Subclass of};
\node[rdfsSubClassOf ,minimum width=82pt  , text=black] at (726.355368355004pt, -461.2573117389028pt)   (property16) {Subclass of};
\node[rdfsSubClassOf ,minimum width=82pt  , text=black] at (1318.9305452273738pt, 45.939166406875984pt)   (property17) {Subclass of};
\node[owlObjectProperty ,minimum width=109pt  , text=black] at (872.5518902451007pt, -262.3232245218421pt)   (property18) {manufactured by};
\node[owlDatatypeProperty ,minimum width=48pt  , text=black] at (272.9803116825276pt, -601.3761639962884pt)   (property19) {isIBA};
\node[owlObjectProperty ,minimum width=83pt  , text=black] at (-142.45219271259987pt, -662.3072233662754pt)   (property20) {ofIngredient};
\node[owlDatatypeProperty ,minimum width=81pt  , text=black] at (388.8890249312997pt, -828.2751456420385pt)   (property21) {preparation};
\node[owlObjectProperty ,minimum width=107pt  , text=black] at (616.837873685871pt, -503.0911442471146pt)   (property22) {hasFlavorProfile};
\node[owlObjectProperty ,minimum width=71pt  , text=black] at (1453.6070119759615pt, -311.56310157930125pt)   (property23) {located in};
\end{tikzpicture}
}
 \end{center}

  \caption{VOWL rendering of \emph{DrinkOnto}. Dotted arrows are
  \texttt{rdfs:subClassOf}.}
  \label{fig:ontology}
\end{figure}
